\documentclass[11pt,a4paper,UTF8,fontset=fandol,oneside,openany]{ctexrep}
\usepackage[margin=2.35cm]{geometry}
\usepackage{amsmath,amssymb,mathtools,bm,booktabs,array,graphicx,hyperref,xcolor,enumitem,fancyhdr}
\usepackage{tikz}
\usepackage[section]{placeins}
\usepackage{float}
\usetikzlibrary{arrows.meta,positioning,calc}
\tikzset{gene/.style={draw,rounded corners,minimum width=12mm,minimum height=8mm,fill=blue!6},act/.style={-{Stealth},thick,blue!65!black},rep/.style={-{Bar[width=6pt]},thick,red!70!black}}
\hypersetup{colorlinks=true,linkcolor=blue!50!black,urlcolor=blue!45!black,citecolor=blue!45!black}
\setlength{\headheight}{15pt}
\setlength{\parindent}{2em}\setlength{\parskip}{.15em}
\setcounter{tocdepth}{1}\numberwithin{equation}{chapter}
\newcommand{\dd}{\mathrm d}\newcommand{\R}{\mathbb R}\newcommand{\ssv}{\mathrm{ss}}
\pagestyle{fancy}\fancyhf{}\fancyhead[L]{最小基因调控回路：从反应到动力学}\fancyhead[R]{\thepage}
\ctexset{chapter={format=\Large\bfseries,beforeskip=0pt,afterskip=18pt},section={format=\large\bfseries\raggedright,beforeskip=2.5ex,afterskip=1ex}}
\fancypagestyle{plain}{\fancyhf{}\fancyhead[L]{最小基因调控回路：从反应到动力学}\fancyhead[R]{\thepage}}
\title{\bfseries 最小基因调控回路的理论基础\\[4pt]\large 从化学反应到一、二、三组分系统及真实网络}
\author{可复现学习笔记（含源代码与数值图）}
\date{2026 年 10 月}
